\section{Methodology: Comparison Criteria}
\label{sec:methodology_criteria}

Candidates are compared using four criteria: stability, interpretability, resistance to gaming, and long-term suitability.
The comparative phase (Phase II) applies these criteria uniformly to M1--M5 before any PoC metric decision.

The scoring labels are local to each analysis step. C2, C3, and C4 each use five criterion-specific dimensions:
\texttt{C2\_D1}--\texttt{C2\_D5}, \texttt{C3\_D1}--\texttt{C3\_D5}, and
\texttt{C4\_D1}--\texttt{C4\_D5}. These are separate from the smart-contract feasibility rubric, which later uses
ten dimensions, D1--D10, in Section~\ref{sec:feasibility_rubric_full}.

\subsection{Evaluation unit and normalization policy}
\label{sec:evaluation_unit}

The primary analysis unit is \texttt{(project\_id, window\_id, developer)} for developer-level metrics and
\texttt{(project\_id, window\_id)} for aggregate metrics. Because window durations and project sizes vary, the pipeline
produces both raw and normalized values:
\begin{itemize}[leftmargin=*]
  \item per-day normalization for cadence-sensitive measures;
  \item per-KLOC normalization where size effects can dominate;
  \item within-project shares (ownership) to control for absolute scale.
\end{itemize}
Windows shorter than seven days are retained for auditability but flagged as \texttt{sensitivity\_only}; Phase II reports
both full-dataset results and robust results excluding those very short windows when duration-sensitive metrics are evaluated.

Automation is handled as an explicit population rule. Rows flagged with \texttt{is\_bot} remain in the clean all-account
tables for auditability and project-activity totals. However, developer-behavior metrics use human-only outputs as the
primary population. In the centralized Phase I dataset, M3 cadence and M4 ownership/concentration are evaluated from
\path{out/metrics_window_dev_human_all.csv} and \path{out/ownership_summary_human_all.csv}. Bot summaries are reported
as sensitivity evidence rather than silently mixed into maintainer-behavior claims.

\subsection{C1 Stability}
\label{sec:criterion_stability}
Stability is evaluated across consecutive version windows. The analysis measures volatility (e.g., dispersion across windows) and
rank stability (e.g., correlation of developer rankings between adjacent windows). Cross-project comparison is performed
after simple normalizations (e.g., per day or per KLOC) to reduce artifacts from project size and release cadence.
The output set comprises dispersion statistics, adjacent-window rank correlations, and sensitivity checks on timestamp semantics.

\subsection{C2 Interpretability}
\label{sec:criterion_interpretability}
Interpretability is assessed via semantic clarity (plain meaning), unit clarity (counts/shares/time), and normalization clarity.
Metrics that require complex tooling or ambiguous definitions receive lower interpretability.

\subsection{C3 Resistance to Gaming}
\label{sec:criterion_gaming}
The analysis examines attack surfaces and manipulation cost for each metric. Since individual metrics can be misleading, the interpretation avoids strong causal claims
and treats gaming risks comparatively (e.g., commit-splitting inflates frequency but becomes detectable when churn remains low).
The threat model follows reputation-system evaluation research that compares metric behavior under hostile-agent scenarios
and other explicit attacks~\cite{koutrouli2015reputation,schlosser2004reputation}. Strong causal interpretations are avoided because
code-complexity measures can be useful as complementary indicators without any single measure having a strong association with
vulnerability~\cite{scvulncomplexity}.

\subsection{C4 Long-Term Suitability}
\label{sec:criterion_longterm}
The assessment considers whether each metric remains informative over long histories and across lifecycle phases, and whether it supports
incremental updates without full recomputation. This lifecycle emphasis follows maintainability research that treats changeability,
evolution, and paradigm-sensitive measurement as distinct software-quality concerns
~\cite{kumar2012maintainability,kapllani2020maintainability,syedmohamad2026maintainability}.

\subsection{Criterion-to-output mapping}
\label{sec:criterion_mapping}

\begin{table}[H]
\centering
\footnotesize
\renewcommand{\arraystretch}{1.15}
\begin{tabularx}{\textwidth}{>{\raggedright\arraybackslash}p{1.8cm}>{\raggedright\arraybackslash}p{5.6cm}>{\raggedright\arraybackslash}X}
\toprule
\textbf{Criterion} & \textbf{Phase II evidence} & \textbf{Primary data sources} \\
\midrule
C1 Stability &
Window-level volatility, adjacent-window rank correlation, robustness to normalization choices. &
\path{out/metrics_window_dev_human_normalized_all.csv}, \path{out/metrics_window_totals_human_normalized_all.csv}, \path{out/window_duration_flags_all.csv} \\
C2 Interpretability &
Structured rubric (semantic clarity, unit clarity, normalization clarity, tooling transparency). &
Metric definitions plus extracted fields in \path{out/} artifacts \\
C3 Gaming resistance &
Attack-surface matrix and manipulation detectability checks under multi-metric reading. &
\path{out/metrics_window_dev_human_all.csv}, \path{out/ownership_dev_shares_human_all.csv}, \path{out/bot_summary_by_project_all.csv} \\
C4 Long-term suitability &
Informativeness over time, lifecycle robustness, and incremental update feasibility. &
\path{out/loc_per_tag_all.csv}, \path{data/windows_all.csv}, feasibility rubric dimensions \\
\bottomrule
\end{tabularx}
\caption{Operational mapping from methodology criteria to evidence sources.}
\label{tab:criterion_mapping}
\end{table}

\subsection{Decision protocol (predefined, not executed in Phase I)}
\label{sec:decision_protocol}

The decision protocol is fixed before Phase II analysis:
\begin{enumerate}[leftmargin=*]
  \item compute comparative evidence for all candidates under C1--C4;
  \item score smart-contract feasibility using the rubric in Section~\ref{sec:feasibility_rubric_full};
  \item eliminate any candidate that violates feasibility exclusion criteria;
  \item select one PoC metric with the best overall trade-off between comparative behavior and feasibility.
\end{enumerate}
No step of this decision protocol is executed in Phase I.





